\documentclass[planificacion.tex]{subfiles}
\begin{document}

\section*{\centering PROGRAMA \ano}

\setstretch{1.8}

\subsection*{Primera Parte:}  	Termodinámica. Calor y Temperatura. Definiciones. Unidades. Escalas termométricas fundamentales, Celsius, Fahrenheit y Kelvin. Conversión de unidades. Transmisión de calor por conducción, por radiación y convección. Formulas. Coeficientes de transmisión. 
Dilatación lineal, superficial y cubica, coeficientes. Calorimetría, definición, unidades, capacidad calorífica, calor especifico, unidades. Ecuación fundamental de la calorimetría. Trabajo termodinámico de expansión y circulación. Representación gráfica de los distintos tipos de trabajos. Definición principio cero de la termodinámica. Sistemas termodinámicos,  tipos, variables fundamentales diagrama P-V. Principio de conservación de la energía. Tipos de transformaciones termodinámicas, isotérmicas, isobáricas, isocóricas, adiabáticas, politrópicas. Graficas. Energía interna y su relación con la entropía. Procesos termodinámicos ideales e irreversibles. Segundo principio de la termodinámica. Relación con los enunciados de  Clausius y Kelvin-Planck. Maquina ideal de Carnot ciclo ideal. Diagramas de procesos entrópicos.
\subsection*{Segunda Parte:}  Maquinas térmicas. Definición, generalidades, usos. Motores de combustión interna. Descripción y funcionamiento.  Ciclos o tiempos termodinámicos. Ciclo Otto, rendimiento real e ideal. Fórmulas. Tiempos o ciclos termodinámicos en los motores de combustión interna. Ciclo Diesel. Rendimiento térmico. Interpretación gráfica de las variables intervinientes en un ciclo ideal de combustión interna. Ciclo semi Diesel. Turbinas de gas, ciclo Brayton. Descripción y funcionamiento.    Diferenciación de rendimientos de los distintos ciclos de combustión interna. Gráficas. Maquinas frigoríficas, clasificación, ciclo. 
Transformaciones de un sistema gaseoso. Fenómenos de vaporización. Tablas de vapor. Entalpia del líquido y del vapor en máquinas refrigerantes. Conocimiento de los Principios de Carnot para ciclos de vapor. Diagramas, título de vapor,  enunciado de Mollier.  Enunciado de William John Rankine. Reconocimiento de los efectos de presión y temperatura en los ciclos de vapor. Conocer los principios vinculados con los ciclos de gas. Funcionamiento de una turbomáquina motora a gas. Propiedades técnico-térmicas de los ciclos de vapor y gas. Ciclos de gas de potencia y de refrigeración según Brayton. Funciones del compresor, turbina y dispositivos acoplados en la generación eléctrica. Eficiencias y rendimientos entre los ciclos combinados.

\end{document}

